Shelfmark 106, scan 14.

\scanpage{14}
\mspage{7})

which we again compose as before. We do indeed have \deleted{\unsure{two|one}} \unsure{functors|functor}
\diagram{\begin{tikzcd}
\deleted{\unsure{C\longrightarrow}}&\widetilde C\arrow[r]&W^{-1}C
\end{tikzcd}}

\deleted{The first of these two functors does indeed transform the elements of $W$ into invertible arrows. \deleted{We have} For
\diagram{\begin{tikzcd}
x\arrow[r,"f","\sim"']&y
\end{tikzcd}}
\unsure{$f\in W$}, we consider
\[\text{let }F=[f]\]
\diagram{\begin{tikzcd}
y\arrow[dr,"1_y"']&x\arrow[d,"f","\sim"']\\
&y
\end{tikzcd}}
\[\text{let }G=[1_y,f].\]
\unsure{And we find the composites} \unsure{and we therefore have}
\[\uline{GF}=\unsure{[\,]^{x}}\]
}
\unsure{let us define the} $=$
\diagram{\begin{tikzcd}C\arrow[r]&\widetilde C\end{tikzcd}}
which is going to be a little more delicate. Let $f:x\longrightarrow y$. We choose a factorization
\diagram{\begin{tikzcd}[column sep=large,row sep=large]
x\arrow[dr,hook,"f_0"']&y\arrow[dr,bend left=30,dotted,"i"]\arrow[dr,dashed,"\beta","\sim"']&\\
&\bar y\arrow[u,"\unsure{\alpha}","\sim"']\arrow[r,dotted,"\unsure{g}","\sim"']&\bar y'
\end{tikzcd}}
\inserted{\unsure{thus $\bar y$ cof}}
and, with the \unsure{w.k lifting lemma} (here $y$ being cofibrant) we can complete with dotted lines, where $\beta$ has a section $i$.

We then consider $[g f_0,i]$, we must show that this does not depend on the choices made. For a choice $[f',i']$
\diagram{\begin{tikzcd}
x\arrow[r,dashed,"f"]\arrow[dr,"f'"']&y\arrow[d,"i'","\sim"',shift right=1]\\
&y'\arrow[u,dashed,"\alpha","\sim"',shift right=1]
\end{tikzcd}}
such that there exists \deleted{\unsure{a|a}} a retraction $\alpha$ from $\widetilde y$ onto $y$, such that
\[\alpha f'=f.\]

\margin{We have reduced to factorizations $f=\unsure{\alpha}f_0$, $\unsure{f_0}\in\unsure{\operatorname{cof}}$ and $\unsure{\alpha}\in W$, such that $\unsure{\alpha}$ has a section.}

Consider two such choices, we must see that they are in a $=$ connected component of $\underline{\operatorname{Ch}}_0(x,y)$. But for two factorizations of an arrow $f:x\longrightarrow y$ as $\alpha f_0$ and $\alpha'f'_0$, $f_0,f'_0\in\operatorname{Cof}$, $\alpha,\alpha'\in W$, \deleted{\unsure{$x\longrightarrow y$}} we can dominate them by a third \deleted{\unsure{$f''$}} $\alpha'',f''_0$ which \deleted{\unsure{\diagram{\begin{tikzcd}x\arrow[r]&y''\end{tikzcd}}}} \deleted{\unsure{makes it possible to}} \inserted{\unsure{receives|reduces}} the two \unsure{as we see in the diagrams}

\margin{This functor is bijective on objects, surjective on arrows, \unsure{as prev.}. We check easily [for this it suffices \unsure{in a product} to see that for ev. $f\in W$, $f:x\longrightarrow y$, $[f]^{-1}$ comes from $\widetilde C$, we take
\diagram{\begin{tikzcd}
x\arrow[r,"\unsure{f}"]\arrow[dr,"\sim"']&y\arrow[d,shift left=1,"\alpha"]\\
&\unsure{\bar y}\arrow[u,very thick,shift left=1,"\unsure{i}","\sim"']
\end{tikzcd}}
\unsure{(the two isos)}]
\diagram{\begin{tikzcd}
y\arrow[dr,"\unsure{1_y}"']&\unsure{x}\arrow[d,"\unsure{f}","\sim"',shift right=1]\\
&y\arrow[u,"\unsure{\alpha}"',shift right=1]
\end{tikzcd}}
}

\figure{\begin{tikzpicture}[x=.4cm,y=.4cm]
\draw (0,.2) .. controls (-1.4,.4) and (-.8,3.2) .. (1.8,5.5) .. controls (3.8,7.3) and (5.4,7.8) .. (5.1,5.4) .. controls (4.8,2.1) and (1.8,-.5) .. (0,.2);
\end{tikzpicture}}
\mspage{13}
